\section{Sujet n\textsuperscript{o}\,14}
\noindent\begin{minipage}{0.5\linewidth}\footnotesize Office du Baccalauréat du Cameroun\end{minipage}%
\hfill\begin{minipage}{0.45\linewidth}\footnotesize\raggedleft Examen : \textbf{Baccalauréat} · Série : \textbf{C/E}\\
Épreuve : \textbf{Mathématiques} · Durée : \textbf{4\,h} · Coef. : \textbf{7}\end{minipage}
\par\smallskip{\color{onbo}\rule{\linewidth}{1pt}}

\subsection*{Partie A — Évaluation des ressources \hfill (15 points)}

\noindent\textbf{Exercice 1.} \hfill\textit{(3 points)}\\
Dans une usine d'embouteillage d'eau minérale de Bafoussam, on a constaté que, sur la
chaîne de production, $5\,\%$ des bouteilles présentent un défaut de bouchage. Le service
qualité prélève au hasard un échantillon de $n=8$ bouteilles dans un lot très important
(le tirage est assimilé à un tirage successif avec remise). On note $X$ la variable
aléatoire égale au nombre de bouteilles défectueuses parmi les $8$ prélevées.
\begin{enumerate}[label=\arabic*)]
\item Justifier que $X$ suit une loi binomiale dont on précisera les paramètres. \hfill\textit{(0,75 pt)}
\item Calculer, à $10^{-4}$ près, la probabilité de l'événement $A$ : « l'échantillon ne
contient aucune bouteille défectueuse ». \hfill\textit{(0,75 pt)}
\item Calculer, à $10^{-4}$ près, la probabilité de l'événement $B$ : « l'échantillon
contient au moins deux bouteilles défectueuses ». \hfill\textit{(1 pt)}
\item Le lot est refusé si l'échantillon contient au moins deux bouteilles défectueuses.
Déterminer l'espérance mathématique $E(X)$ et interpréter ce résultat. \hfill\textit{(0,5 pt)}
\end{enumerate}

\smallskip
\noindent\textbf{Exercice 2.} \hfill\textit{(3 points)}\\
Chaque bouteille produite reçoit un numéro de série. Pour la traçabilité, on s'intéresse
à des questions de congruences.
\begin{enumerate}[label=\arabic*)]
\item Déterminer le reste de la division euclidienne de $7^{2025}$ par $11$. \hfill\textit{(1 pt)}
\item Résoudre dans $\Z^2$ l'équation $(E):\ 13x-9y=1$. \hfill\textit{(1 pt)}
\item En déduire l'ensemble des entiers relatifs $x$ tels que $13x\equiv 1\ [9]$, puis
le plus petit entier naturel $x_0$ vérifiant cette congruence. \hfill\textit{(1 pt)}
\end{enumerate}

\smallskip
\noindent\textbf{Exercice 3.} \hfill\textit{(4 points)}\\
Le volume d'eau (en milliers de litres) présent chaque jour dans une cuve de l'usine est
modélisé par la suite $(u_n)$ définie par $u_0=1$ et, pour tout $n\in\N$,
$u_{n+1}=f(u_n)$ où $f$ est la fonction définie sur $[0,+\infty[$ par
$f(x)=\dfrac{4x+2}{x+3}$.
\begin{enumerate}[label=\arabic*)]
\item Étudier les variations de $f$ sur $[0,+\infty[$ et montrer que $f([0,2])\subset[0,2]$. \hfill\textit{(1 pt)}
\item Montrer par récurrence que, pour tout $n\in\N$, $0\le u_n\le u_{n+1}\le 2$. \hfill\textit{(1,25 pt)}
\item Justifier que la suite $(u_n)$ converge et déterminer sa limite $\ell$. \hfill\textit{(1 pt)}
\item On pose $v_n=\dfrac{u_n-2}{u_n+1}$. Montrer que $(v_n)$ est géométrique, puis
exprimer $u_n$ en fonction de $n$. \hfill\textit{(0,75 pt)}
\end{enumerate}

\smallskip
\noindent\textbf{Exercice 4.} \hfill\textit{(5 points)}\\
Le plan est muni d'un repère orthonormé direct $(O\,;\,\vec\imath,\vec\jmath)$ (unité :
le mètre). On modélise l'aire de stockage par ce plan. On considère les points
$A(-1\,;\,2)$, $B(5\,;\,4)$ et $C(3\,;\,-2)$.
\begin{enumerate}[label=\arabic*)]
\item Calculer le produit scalaire $\vec{AB}\cdot\vec{AC}$ et en déduire une valeur
approchée au degré près de l'angle $\widehat{BAC}$. \hfill\textit{(1 pt)}
\item Déterminer l'ensemble $(\Gamma_1)$ des points $M$ du plan tels que
$\vec{MA}\cdot\vec{MB}=0$. Préciser sa nature et ses éléments caractéristiques. \hfill\textit{(1,25 pt)}
\item Déterminer et construire l'ensemble $(\Gamma_2)$ des points $M$ du plan tels que
$MA^2+MB^2=40$. \hfill\textit{(1,5 pt)}
\item Déterminer l'ensemble $(\Gamma_3)$ des points $M$ tels que
$\vec{MA}\cdot\vec{MB}=12$, et préciser sa nature. \hfill\textit{(1,25 pt)}
\end{enumerate}

\subsection*{Partie B — Évaluation des compétences \hfill (5 points)}

\noindent\textit{Le directeur de l'usine d'embouteillage de Bafoussam doit prendre trois
décisions et te sollicite, en tant qu'élève de Terminale~C. \textbf{(i)} Il fait fabriquer
des caisses parallélépipédiques sans couvercle, de base carrée, destinées au transport
des bouteilles : il dispose pour chaque caisse de \SI{12}{\metre\squared} de carton et
veut que le volume de chaque caisse soit \emph{maximal}. \textbf{(ii)} La production
mensuelle, de \num{50000} bouteilles le premier mois, augmente de \SI{6}{\percent} chaque
mois ; il souhaite savoir quand la production mensuelle dépassera \num{100000} bouteilles.
\textbf{(iii)} À la sortie de chaîne, $3\,\%$ des bouteilles sont mal étiquetées ; un
client en prélève $10$ au hasard et refuse la livraison si au moins une bouteille est mal
étiquetée.}

\begin{enumerate}[label=\textbf{Tâche \arabic*.},leftmargin=2.4em]
\item Déterminer les dimensions de la caisse de volume \emph{maximal} et la valeur de ce
volume. \hfill\textit{(1,5 pt)}
\item Déterminer le rang du premier mois à partir duquel la production mensuelle dépassera
\num{100000} bouteilles. \hfill\textit{(1,5 pt)}
\item Déterminer, à $10^{-2}$ près, la probabilité que le client refuse la livraison. \hfill\textit{(1,5 pt)}
\end{enumerate}
\par\smallskip{\footnotesize\itshape Présentation générale : 0,5 point.}

% ════════════════════════════════════════════════════
\corrige{du sujet 14}

\noindent\textbf{Partie A — Exercice 1.}
1) On répète $n=8$ fois, de façon indépendante (tirage avec remise), une épreuve de
Bernoulli dont le succès « bouteille défectueuse » a pour probabilité $p=0{,}05$. Donc
$X\hookrightarrow\mathcal{B}(8\,;\,0{,}05)$.
2) $P(A)=P(X{=}0)=(0{,}95)^8\approx\mathbf{0{,}6634}$.
3) $P(X{=}1)=\binom{8}{1}(0{,}05)(0{,}95)^7\approx0{,}2793$, donc
$P(B)=P(X\ge2)=1-P(X{=}0)-P(X{=}1)\approx1-0{,}6634-0{,}2793=\mathbf{0{,}0573}$.
4) $E(X)=np=8\times0{,}05=0{,}4$ : en moyenne, $0{,}4$ bouteille défectueuse par
échantillon de $8$.

\smallskip
\noindent\textbf{Exercice 2.}
1) Modulo $11$ : $7^2=49\equiv5$, $7^3\equiv35\equiv2$, $7^4\equiv14\equiv3$,
$7^5\equiv21\equiv10\equiv-1$, donc $7^{10}\equiv1\ [11]$. Comme $2025=10\times202+5$,
$7^{2025}\equiv(7^{10})^{202}\cdot7^5\equiv7^5\equiv-1\equiv\mathbf{10}\ [11]$. Le reste est $10$.
2) $13\wedge9=1$ : par l'algorithme d'Euclide, $13=1\times9+4$, $9=2\times4+1$, d'où
$1=9-2\times4=9-2(13-9)=3\times9-2\times13$, soit $13\times(-2)-9\times(-3)=1$. Une
solution particulière est $(-2\,;\,-3)$. Par différence $13(x+2)=9(y+3)$ ; comme
$13\wedge9=1$, Gauss donne $9\mid x+2$, d'où $x=-2+9k$, $y=-3+13k$, $k\in\Z$.
3) $13x\equiv1\ [9]\iff 9\mid 13x-1\iff \exists y,\ 13x-9y=1$, donc $x\equiv-2\equiv7\ [9]$.
L'ensemble est $\{7+9k,\ k\in\Z\}$ ; le plus petit entier naturel est $x_0=\mathbf{7}$
(on vérifie : $13\times7=91=9\times10+1$).

\smallskip
\noindent\textbf{Exercice 3.}
1) $f(x)=\dfrac{4x+2}{x+3}$ ; $f'(x)=\dfrac{4(x+3)-(4x+2)}{(x+3)^2}=\dfrac{10}{(x+3)^2}>0$ :
$f$ est strictement croissante sur $[0,+\infty[$. Sur $[0,2]$, $f$ croissante donne
$f([0,2])=[f(0),f(2)]=[\tfrac23,2]\subset[0,2]$.
2) Initialisation : $u_0=1$, $u_1=f(1)=\tfrac{6}{4}=1{,}5$, donc $0\le u_0\le u_1\le2$.
Hérédité : supposons $0\le u_n\le u_{n+1}\le2$. $f$ étant croissante,
$f(0)\le f(u_n)\le f(u_{n+1})\le f(2)$, soit $\tfrac23\le u_{n+1}\le u_{n+2}\le2$,
d'où $0\le u_{n+1}\le u_{n+2}\le2$. La propriété est héréditaire, donc vraie pour tout $n$.
La suite $(u_n)$ est donc croissante et majorée par $2$.
3) $(u_n)$ est croissante et majorée par $2$, donc elle converge vers une limite
$\ell\in[0,2]$. Par continuité de $f$, $\ell=f(\ell)$, soit $\ell(\ell+3)=4\ell+2$,
$\ell^2-\ell-2=0$, $(\ell-2)(\ell+1)=0$, d'où $\ell=2$ (la racine $-1$ est exclue).
Donc $\boxed{\ell=2}$.
4) $v_n=\dfrac{u_n-2}{u_n+1}$. On calcule
$v_{n+1}=\dfrac{f(u_n)-2}{f(u_n)+1}
=\dfrac{\frac{4u_n+2}{u_n+3}-2}{\frac{4u_n+2}{u_n+3}+1}
=\dfrac{4u_n+2-2(u_n+3)}{4u_n+2+(u_n+3)}
=\dfrac{2u_n-4}{5u_n+5}=\dfrac{2(u_n-2)}{5(u_n+1)}=\dfrac{2}{5}\,v_n$.
Donc $(v_n)$ est géométrique de raison $q=\dfrac25$ et de premier terme
$v_0=\dfrac{1-2}{1+1}=-\dfrac12$. Ainsi $v_n=-\dfrac12\left(\dfrac25\right)^n$. De
$v_n=\dfrac{u_n-2}{u_n+1}$ on tire $u_n=\dfrac{2+v_n}{1-v_n}=\dfrac{4-\left(\frac25\right)^n}{2+\left(\frac25\right)^n}$,
et l'on retrouve $u_n\to2$ quand $n\to+\infty$.

\smallskip
\noindent\textbf{Exercice 4.}
$\vec{AB}=(6\,;\,2)$, $\vec{AC}=(4\,;\,-4)$.
1) $\vec{AB}\cdot\vec{AC}=6\times4+2\times(-4)=24-8=16$. $\|\vec{AB}\|=\sqrt{40}=2\sqrt{10}$,
$\|\vec{AC}\|=\sqrt{32}=4\sqrt2$. $\cos\widehat{BAC}=\dfrac{16}{2\sqrt{10}\cdot4\sqrt2}
=\dfrac{16}{8\sqrt{20}}=\dfrac{2}{2\sqrt5}=\dfrac{1}{\sqrt5}\approx0{,}447$, d'où
$\widehat{BAC}\approx\mathbf{63^\circ}$.
2) $\vec{MA}\cdot\vec{MB}=0$ : $(\Gamma_1)$ est le cercle de diamètre $[AB]$, de centre
$I\big(\tfrac{-1+5}{2}\,;\,\tfrac{2+4}{2}\big)=I(2\,;\,3)$ et de rayon
$\tfrac12 AB=\tfrac12\cdot2\sqrt{10}=\sqrt{10}$.
3) Avec $G$ milieu de $[AB]$ (donc $G=I(2\,;\,3)$), la formule de la médiane donne
$MA^2+MB^2=2\,MI^2+\tfrac12 AB^2$. Or $AB^2=40$, donc $MA^2+MB^2=2\,MI^2+20$. L'équation
$MA^2+MB^2=40$ équivaut à $2\,MI^2=20$, soit $MI^2=10$, $MI=\sqrt{10}$. $(\Gamma_2)$ est le
cercle de centre $I(2\,;\,3)$ et de rayon $\sqrt{10}$ (identique au cercle de diamètre $[AB]$).
4) $\vec{MA}\cdot\vec{MB}=MI^2-\dfrac{AB^2}{4}=MI^2-10$. L'équation
$\vec{MA}\cdot\vec{MB}=12$ équivaut à $MI^2=22$, soit $MI=\sqrt{22}$. $(\Gamma_3)$ est le
cercle de centre $I(2\,;\,3)$ et de rayon $\sqrt{22}$.

\begin{center}
\begin{tikzpicture}[scale=0.5,>=Stealth]
\draw[->,onbmuted] (-5,0)--(9,0) node[right]{$x$};
\draw[->,onbmuted] (0,-5)--(0,8) node[above]{$y$};
\foreach \x in {-4,-2,2,4,6,8} \draw[onbline] (\x,-0.12)--(\x,0.12);
\foreach \y in {-4,-2,2,4,6} \draw[onbline] (-0.12,\y)--(0.12,\y);
\draw[onbo,thick] (2,3) circle (3.1623);
\fill[onbink] (-1,2) circle (3pt) node[below left]{$A$};
\fill[onbink] (5,4) circle (3pt) node[above right]{$B$};
\fill[onbink] (3,-2) circle (3pt) node[below right]{$C$};
\fill[onbo] (2,3) circle (2.5pt) node[above]{$I$};
\end{tikzpicture}
\end{center}
\fig{Le cercle $(\Gamma_2)$ de centre $I(2\,;\,3)$ et de rayon $\sqrt{10}$.}

\smallskip
\noindent\textbf{Partie B — Situation.}
\emph{Tâche 1.} Caisse à base carrée de côté $x$ et de hauteur $h$, sans couvercle.
Aire de carton : base $+$ $4$ faces latérales : $x^2+4xh=12$, d'où $h=\dfrac{12-x^2}{4x}$
($0<x<2\sqrt3$). Volume : $V(x)=x^2h=x^2\cdot\dfrac{12-x^2}{4x}=\dfrac{12x-x^3}{4}=3x-\dfrac{x^3}{4}$.
$V'(x)=3-\dfrac{3x^2}{4}=\dfrac34(4-x^2)$, qui s'annule en $x=2$ (positif) avec $V'>0$ sur
$]0,2[$ et $V'<0$ sur $]2,2\sqrt3[$ : maximum en $x=2$. Alors $h=\dfrac{12-4}{8}=1$ et
$V(2)=3\times2-\dfrac{8}{4}=6-2=\SI{4}{\metre\cubed}$. \textbf{Caisse} : base
$2\,\text{m}\times2\,\text{m}$, hauteur $1\,\text{m}$, volume \textbf{$4\,\text{m}^3$}.

\emph{Tâche 2.} Production du mois $n$ : $P_n=50000\times1{,}06^{\,n-1}$. On résout
$P_n>100000$, soit $1{,}06^{\,n-1}>2$, $n-1>\dfrac{\ln 2}{\ln 1{,}06}\approx\dfrac{0{,}6931}{0{,}0583}\approx11{,}9$,
d'où $n-1\ge12$, $n\ge13$. À partir du \textbf{13\textsuperscript{e} mois}
($P_{13}=50000\times1{,}06^{12}\approx\num{100610}$ bouteilles).

\emph{Tâche 3.} Soit $Y$ le nombre de bouteilles mal étiquetées parmi $10$ :
$Y\hookrightarrow\mathcal{B}(10\,;\,0{,}03)$. Le client refuse si $Y\ge1$ :
$P(Y\ge1)=1-(0{,}97)^{10}\approx1-0{,}7374=\mathbf{0{,}26}$, soit environ \SI{26}{\percent}.
